% !PRINTED 261
\weberpage{261}

If $m$ is a divisor of $p^2-1$, and

\[
\varrho = \gamma^{\frac{p^2-1}{m}},
\]
then $\varrho^m=1$ and there is no lower power of $\varrho$ than the $m$th
that becomes equal to 1. We therefore call $\varrho$ a
primitive $m$th root of unity (in $\mathfrak{E}_{2,p}$). All other $m$th
roots of unity are then powers $\varrho^s$ of $\varrho$, and among these
those and only those are primitive for which $s$ is relatively prime to $m$.

This follows from the fact that every number $\xi$ different from zero
is representable in the form $\xi=\gamma^r$, and that $\xi^m=\gamma^{rm}$ is
$=1$ if and only if $rm$ is divisible by $p^2-1$.

Every number different from zero in $\mathfrak{E}_{2,p}$ is a root of unity
 of degree $p^2-1$.

\begin{center}
\S.\,68.

\gesperrt{The real linear congruence group $L_p$.}
\end{center}

We now proceed, equipped with these means, to the
investigation of the real congruence group $L_p$, which consists of all
substitutions of the form

\[
A = \begin{pmatrix} a, & b \\ c, & d \end{pmatrix}
\]
in which $a, b, c, d$ are real integers taken with respect to the modulus $p$,
which satisfy the condition $ad-bc\equiv 1$.
The group $L_p$ is, according to the earlier notation, to be denoted
by $E_{1,p}$, and its degree, if we exclude the case $p=2$, is:

\[
\frac{p(p^2-1)}{2}.
\]

If $A, U$ are any two elements of a group, we have
already earlier called

\[
U^{-1}AU = A'
\]
the element transformed from $A$ by $U$. $A'^2, A'^3, \ldots$
are transformed from $A^2, A^3, \ldots$, from which it follows that all
elements obtained from one another by transformation have the same
degree. If we take the elements $A$ from any
group $G$, then the transformed elements form a group

\[
U^{-1}GU = G',
\]
isomorphic with $G$, which we call the transformed group of $G$.
